\documentclass[12pt,a4paper,openany]{ctexbook}

\usepackage[top=2.54cm, bottom=2.54cm, left=3.17cm, right=3.17cm]{geometry}
\usepackage{setspace}
\usepackage{array}
\usepackage{tabularx}
\usepackage{xcolor}
\usepackage{fancyhdr}
\usepackage{titlesec}
\usepackage{hyperref}
\usepackage{enumitem}
\usepackage{tikz}
\usetikzlibrary{arrows.meta, positioning, shapes.geometric}

% ========== 页眉页脚 ==========
\pagestyle{fancy}
\fancyhf{}
\fancyhead[C]{\small ESP32-S3 AI 语音对话 —— 项目开发说明}
\fancyfoot[C]{\thepage}
\renewcommand{\headrulewidth}{0.4pt}

% ========== 标题格式 ==========
\titleformat{\chapter}{\centering\heiti\LARGE}{}{0em}{}
\titleformat{\section}{\heiti\large}{\thesection}{1em}{}
\titleformat{\subsection}{\heiti\small}{\thesubsection}{1em}{}
\titlespacing*{\chapter}{0pt}{10pt}{15pt}

% ========== 自定义颜色 ==========
\definecolor{primary}{RGB}{26,60,110}
\definecolor{accent}{RGB}{200,50,50}
\definecolor{success}{RGB}{30,130,70}
\definecolor{warning}{RGB}{200,150,30}
\definecolor{graytext}{RGB}{100,100,100}

% ========== 任务状态命令 ==========
\newcommand{\taskdone}{{\color{success}$\blacksquare$}}
\newcommand{\taskwarn}{{\color{warning}$\blacksquare$}}
\newcommand{\tasktodo}{{\color{graytext}$\blacksquare$}}
\newcommand{\taskfail}{{\color{accent}$\blacksquare$}}

\begin{document}

% ==========================================
% 封面
% ==========================================
\begin{titlepage}
    \vspace*{2cm}
    \begin{center}
        {\heiti\Huge \color{primary} ESP32-S3 AI 语音对话}\\[8pt]
        {\kaishu\LARGE 项目开发说明}\\[20pt]

        \vfill

        \begin{tikzpicture}[scale=1.2]
            \draw[thick,primary,rounded corners] (-5.5,-0.5) rectangle (5.5,0.5);
            \node[primary] at (0,0) {\kaishu\Large INMP441 $\rightarrow$ 百度 ASR $\rightarrow$ DeepSeek $\rightarrow$ 示波器显示};
        \end{tikzpicture}

        \vspace{30pt}

        {\large
         \begin{tabular}{rl}
            \heiti 硬件平台： & \songti ESP32-S3 (RYMCU-DevKitC N8R2) \\
            \heiti 麦克风：   & \songti INMP441 (I2S 标准模式) \\
            \heiti 开发环境： & \songti PlatformIO (VS Code) + Keil MDK \\
            \heiti 文档版本： & \songti v2.1
         \end{tabular}
        }

        \vspace{20pt}
    \end{center}
\end{titlepage}

% ==========================================
% 目录
% ==========================================
\tableofcontents
\thispagestyle{empty}
\newpage

% ==========================================
% 一、项目概述
% ==========================================
\chapter{项目概述}

\section{项目目标}

通过 INMP441 采集语音，调用百度语音识别 (ASR) 和 DeepSeek 大语言模型 (LLM) 生成回复，最终通过 UART 将英文回复发送给 STM32 主机，在示波器 (XY 模式) 上显示文字。

\section{当前阶段}

当前项目处于 \textbf{第一阶段}（ESP32 端适配），已完成 ESP32 上录音、ASR、LLM 全链路的开发和验证，正进行示波器显示方向的适配改造。

\section{快速开始}

\subsection*{编译与上传}
\begin{verbatim}
# 进入项目目录
cd AI_Chat

# 编译
pio run

# 上传（ESP32-S3 默认使用 COM9 USB-Serial/JTAG 端口）
pio run -t upload --upload-port COM9

# 串口监视器
pio device monitor -b 115200 --port COM9
\end{verbatim}

\subsection*{串口命令}
\begin{center}
\begin{tabularx}{0.7\textwidth}{|c|X|}
    \hline
    \heiti 命令 & \heiti 功能\\
    \hline
    \texttt{s / S} & 开始一轮对话（录音→ASR→LLM→回复）\\
    \hline
    \texttt{w} & 重新连接 WiFi\\
    \hline
    \texttt{h / H} & 打印帮助信息\\
    \hline
\end{tabularx}
\end{center}

\subsection*{端口说明}
\begin{itemize}
    \item ESP32-S3 使用 USB-Serial/JTAG 内置串口，在 Windows 上通常识别为 \texttt{COM9}。
    \item 如果上次上传后未断电，再次上传可能端口被占用，此时需按一下板子上的 \texttt{RST} 按钮再试。
    \item 蓝牙虚拟 COM 口（如 COM3/COM5/COM6/COM7/COM8/COM13/COM14/COM18/COM19）不可用于 ESP32 上传。
\end{itemize}

\section{系统架构}

\begin{center}
\vspace{10pt}
\begin{tikzpicture}[
    node distance=1.8cm,
    box/.style={rectangle, draw, rounded corners, thick, align=center, text width=2.5cm, minimum height=1.0cm, font=\small},
    arrow/.style={-{Latex[length=2.5mm]}, thick}
]

\node[box, fill=blue!10!white, draw=blue] (mic) {INMP441 \\ I2S 麦克风};
\node[box, fill=red!10!white, draw=red, right=of mic] (esp) {ESP32-S3};
\node[box, fill=green!10!white, draw=green, right=of esp, yshift=1.2cm] (baidu) {百度 ASR};
\node[box, fill=green!10!white, draw=green, right=of esp, yshift=-1.2cm] (deepseek) {DeepSeek LLM};
\node[box, fill=orange!10!white, draw=orange, below=of esp] (stm) {STM32F407};
\node[box, fill=purple!10!white, draw=purple, below=of stm] (scope) {示波器 XY 显示};

\draw[arrow] (mic) -- node[above,font=\tiny] {I2S 32bit $\to$ 16bit PCM} (esp);
\draw[arrow] (esp) -- node[right,font=\tiny] {HTTP POST 裸PCM} (baidu);
\draw[arrow] (esp) -- node[right,font=\tiny] {HTTP POST JSON} (deepseek);
\draw[arrow] (baidu) -- node[left,font=\tiny] {识别文本} (esp);
\draw[arrow] (deepseek) -- node[left,font=\tiny] {回复文本} (esp);
\draw[arrow] (esp) -- node[right,font=\tiny] {UART 115200} (stm);
\draw[arrow] (stm) -- node[right,font=\tiny] {DAC X/Y} (scope);

\end{tikzpicture}
\vspace{10pt}
\end{center}

% ==========================================
% 二、功能模块状态
% ==========================================
\chapter{功能模块状态}

\begin{center}
\renewcommand{\arraystretch}{1.3}
\begin{tabularx}{\textwidth}{|c|X|c|X|}
    \hline
    \heiti 模块 & \heiti 功能 & \heiti 状态 & \heiti 注意点\\
    \hline
    音频采集 & INMP441 I2S 驱动 (16kHz, 16bit, 单声道) & \taskdone 已完成 & 标准 I2S RX，不能用 PDM 模式 \\
    \hline
    WiFi 联网 & Station 模式连手机热点 & \taskdone 已完成 & 无自动重连，需手动 'w' 命令 \\
    \hline
    百度 ASR & 裸 PCM POST 语音识别 & \taskdone 已完成 & dev\_pid=1537，token 缓存 29 天 \\
    \hline
    DeepSeek LLM & HTTP POST 对话 & \taskdone 已完成 & 无 /v1/ 前缀，HTTP/1.0 强制，模型 deepseek-v4-flash \\
    \hline
    全链路串通 & 录音$\to$ASR$\to$LLM$\to$输出 & \taskdone 已完成 & 已验证输出故事，中文偶尔乱码 \\
    \hline
    内存管理 & PSRAM 利用 & \taskfail 有问题 & PSRAM ID 读取错误，fallback 到 malloc \\
    \hline
    稳定性 & 防崩溃 & \taskfail 有问题 & 5秒录音+HTTP 可能触发 StoreProhibited \\
    \hline
    DeepSeek 提示词 & 改为英文回复 & \taskwarn 待修改 & 适配示波器 ASCII 显示需要 \\
    \hline
    UART 通信 & ESP32 $\to$ STM32 & \tasktodo 未开始 & 需硬件连接 PA9/PA10 \\
    \hline
    STM32 接收 & 接收文本并显示 & \tasktodo 未开始 & FreeRTOS 任务轮询 USART1 \\
    \hline
    示波器显示 & Terminal\_Print 调用 & \tasktodo 未开始 & 矢量字符仅支持 ASCII \\
    \hline
\end{tabularx}
\end{center}

\vspace{10pt}

\textbf{图例：}
\quad {\color{success}$\blacksquare$} 已完成 \qquad
{\color{warning}$\blacksquare$} 待修改 \qquad
{\color{accent}$\blacksquare$} 有问题 \qquad
{\color{graytext}$\blacksquare$} 未开始

% ==========================================
% 三、已完成模块详细记录
% ==========================================
\chapter{已完成模块详细记录}

包含关键实现细节和"坑"，方便后续排查。

\section{麦克风驱动 INMP441}

\subsection*{关键文件}
\begin{itemize}
    \item \texttt{lib/microphone/src/microphone.cpp} —— I2S 驱动封装
    \item \texttt{include/config.h} —— 引脚定义
\end{itemize}

\subsection*{配置参数}
\begin{verbatim}
I2S 模式:      I2S_MODE_MASTER | I2S_MODE_RX
采样率:        16000 Hz
位深:          I2S_BITS_PER_SAMPLE_32BIT
通道格式:      I2S_CHANNEL_FMT_ONLY_LEFT
通信格式:      I2S_COMM_FORMAT_STAND_I2S
DMA:           8 缓冲 × 256 样本
引脚:          SCK=GPIO15, CLK=GPIO5, DATA=GPIO4, L/R=GND
\end{verbatim}

\subsection*{注意点}
\begin{itemize}
    \item \taskfail \textbf{不能用 \texttt{I2S\_MODE\_PDM}}！INMP441 不是 PDM 数字麦克风，它输出标准 I2S 信号。误用 PDM 模式会读出全 0x8729 噪声。
    \item 读取 32-bit 样本后右移 16 位得到 16-bit PCM：\texttt{(int16\_t)(temp[i] >> 16)}。
    \item 采样率 16kHz，10 秒录音需 320KB 缓冲区。
    \item \taskfail \textbf{废弃旧驱动} \texttt{lib/inmp441/} 使用了 \texttt{I2S\_MODE\_PDM}，不可用于此项目。当前使用的正确驱动是 \texttt{lib/microphone/src/}（标准 I2S RX 模式）。
\end{itemize}

\section{WiFi 联网}

\subsection*{配置}
\begin{verbatim}
SSID:     Xiaomi 15 Pro
Password: you3211231234
模式:     WIFI_STA
超时:     40次 × 500ms = 20秒
\end{verbatim}

\subsection*{注意点}
\begin{itemize}
    \item 连接成功会打印 IP 地址（如 10.189.159.90）。
    \item 无断线自动重连，断线后需手动输入 \texttt{w} 命令。
\end{itemize}

\section{百度语音识别 ASR}

\subsection*{接口细节}
\begin{description}
    \item[Token 获取] POST \texttt{https://aip.baidubce.com/oauth/2.0/token}，缓存 29 天。
    \item[语音识别] POST \texttt{https://vop.baidu.com/server\_api}，Content-Type: \texttt{audio/pcm;rate=16000}，dev\_pid=1537。
    \item[数据格式] \textbf{裸 PCM 二进制}，不要 Base64。
    \item[返回] 识别文本在 \texttt{result[0]} 中。
\end{description}

\section{DeepSeek 大语言模型}

\subsection*{接口细节}
\begin{verbatim}
端点:     https://api.deepseek.com/chat/completions
模型:     deepseek-v4-flash
HTTP:     HTTP/1.0 (useHTTP10(true))
超时:     连接10s, HTTP 30s
\end{verbatim}

\subsection*{注意点}
\begin{itemize}
    \item \taskfail \textbf{地址不含 \texttt{/v1/}}。旧版 \texttt{/v1/chat/completions} 已迁移失效。
    \item \texttt{deepseek-chat} 模型 2026/07/24 废弃，当前用 \texttt{deepseek-v4-flash}。
    \item \taskfail \textbf{必须开 HTTP/1.0}，否则 \texttt{WiFiClientSecure::getString()} 在 chunked 编码时报错 \texttt{UNKNOWN ERROR CODE 004C}。
\end{itemize}

\section{全链路对话流程}

\begin{enumerate}
    \item 上电初始化：麦克风 $\rightarrow$ WiFi $\rightarrow$ 打印提示
    \item 串口输入 \texttt{s} 触发录音
    \item 录音中可再发 \texttt{s} 提前停止
    \item 自动执行：录音 $\rightarrow$ 百度 ASR $\rightarrow$ DeepSeek $\rightarrow$ 串口输出
\end{enumerate}

已验证输出故事（"月光下的约定"、"最后一盏灯"等），中文回复偶尔乱码，改为英文后自然解决。

% ==========================================
% 四、待解决问题
% ==========================================
\chapter{待解决问题}

\section{PSRAM 无法使用}

\subsection*{现象}
\begin{verbatim}
E (208) psram: PSRAM ID read error: 0x00ffffff, PSRAM chip not found
⚠️ PSRAM 未启用，录音时长受限
\end{verbatim}

\subsection*{分析}
板载标注 2MB PSRAM，Boot log 也显示 ``Embedded PSRAM 8MB''，但初始化失败。尝试 \texttt{board\_build.psram\_mode = opi} 导致 boot crash（该板 Flash 是 Quad 模式）。

\subsection*{当前状态}
自动 fallback 到 \texttt{malloc()}。\texttt{config.h} 中 \texttt{MAX\_RECORD\_SEC=10} 写死，\texttt{start\_ai\_dialog()} 内检测 PSRAM 失败后降级到 5 秒上限。实测 5 秒录音(160KB) + HTTP 请求可能堆溢出崩溃。

\subsection*{注意点}
\begin{itemize}
    \item 部分 N8R2 批次可能未焊接 PSRAM，需确认硬件。
    \item 短期内将录音限至 3 秒可避免崩溃，但需修改代码硬限。
\end{itemize}

\section{StoreProhibited 崩溃}

\subsection*{触发条件}
录音 5 秒 + \texttt{baidu\_asr()} HTTP POST 时偶发。

\subsection*{原因}
\texttt{malloc(160KB)} 接近堆上限，HTTP 再申请内存时越界。

\subsection*{解决办法}
录音上限降至 3 秒（约 96KB），预留 HTTP 所需内存。

% ==========================================
% 五、阶段划分与任务清单
% ==========================================
\chapter{阶段划分与任务清单}

% ---------------------------
\section{第一阶段：ESP32 端适配（当前阶段）}

\subsection*{当前阶段已完成}
\begin{itemize}
    \item[\taskdone] INMP441 I2S 驱动 (16kHz, 16bit, 单声道) — 使用 \texttt{lib/microphone/} 标准 I2S RX
    \item[\taskdone] WiFi 联网
    \item[\taskdone] 百度 ASR 集成 — 裸 PCM POST，dev\_pid=1537
    \item[\taskdone] DeepSeek 集成 — HTTP/1.0，模型 deepseek-v4-flash
    \item[\taskdone] 全链路：录音 $\to$ ASR $\to$ LLM $\to$ 输出
\end{itemize}

\subsection*{当前阶段未完成}
\begin{itemize}
    \item[\taskwarn] \textbf{修改 DeepSeek system prompt 为英文}
    \begin{itemize}
        \item 当前代码中 system prompt 为中文（\texttt{"你是一个有帮助的AI助手。请用简洁的中文回答用户的问题。"}）
        \item 需改为纯英文 + ASCII（无中文、Emoji、特殊符号），限制输出长度（示波器 10--15 行）
        \item 参见 \texttt{src/main.cpp} 中 \texttt{deepseek\_chat()} 函数的 \texttt{sys\_msg["content"]}
    \end{itemize}

    \item[\taskfail] \textbf{修复录音稳定性}
    \begin{itemize}
        \item \texttt{config.h} 中 \texttt{MAX\_RECORD\_SEC=10}（写死10秒），PSRAM 不可用时 fallback 到 5 秒
        \item 需要：将上限硬限制为 3 秒（约 96KB），避免 \texttt{malloc(160KB)} 堆溢出触发 StoreProhibited
        \item 当前 \texttt{start\_ai\_dialog()} 中已有 PSRAM fallback 逻辑但未硬限 3 秒
    \end{itemize}

    \item[\tasktodo] \textbf{添加 UART 输出}
    \begin{itemize}
        \item 使用 \texttt{HardwareSerial}（如 Serial1）初始化 UART，避免与 USB Serial 冲突
        \item DeepSeek 回复同时输出到 USB Serial（调试）和 UART（送 STM32）
        \item 波特率 115200，引脚避开 I2S 占用的 GPIO4/5/15
        \item API：\texttt{Serial1.begin(115200, SERIAL\_8N1, RX\_PIN, TX\_PIN)}
    \end{itemize}
\end{itemize}

\subsection*{注意点}
\begin{itemize}
    \item PSRAM 不可用（ID 读取错误 0x00ffffff），所有 malloc 用普通堆内存。
    \item 录音超过 3 秒（96KB）可能触发 StoreProhibited 崩溃，第一阶段必须限时。
    \item UART 引脚避免与 I2S 冲突（I2S 用了 GPIO4/5/15），建议用 GPIO16(TX)/GPIO17(RX) 或 ESP32-S3 原生 GPIO43(TX)/GPIO44(RX)。
    \item 修改 system prompt 为英文后，中文乱码问题自然消失（无需额外处理）。
    \item ESP32-S3 上传端口固定为 COM9（USB-Serial/JTAG），蓝牙虚拟 COM 口不可用于上传。
\end{itemize}

% ---------------------------
\section{第二阶段：STM32 端开发}

\subsection*{任务}
\begin{enumerate}[leftmargin=2em, label=\textbf{\arabic*.}]
    \item \textbf{确认 USART1 硬件连接}
    \begin{itemize}
        \item ESP32 TX $\to$ STM32 PA10 (USART1 RX)
        \item 共地
        \item 波特率 115200, 8N1
    \end{itemize}

    \item \textbf{编写 UART 接收任务}
    \begin{itemize}
        \item FreeRTOS 任务 \texttt{vUartRxTask}
        \item \texttt{HAL\_UART\_Receive\_IT()} 或 DMA 接收
        \item 环形缓冲区存储字符
    \end{itemize}

    \item \textbf{文本分行与显示}
    \begin{itemize}
        \item 遇到 \texttt{\textbackslash n} 调用 \texttt{DRAW\_Terminal\_Print()}
        \item 设置字号和行距
        \item 长文本自动分页
    \end{itemize}
\end{enumerate}

\subsection*{注意点}
\begin{itemize}
    \item 矢量字符只支持 ASCII（A-Z, a-z, 0-9, 标点），不支持中文。
    \item PA9/PA10 重映射需确认 GPIO\_AF7\_USART1。
    \item 示波器 XY 模式需调高余晖/持久时间防止闪烁。
\end{itemize}

% ---------------------------
\section{第三阶段：系统联调}

\subsection*{任务}
\begin{enumerate}[leftmargin=2em, label=\textbf{\arabic*.}]
    \item 上下位机联合测试
    \item 调整示波器显示效果（亮度、字号、位置）
    \item 优化端到端延迟
    \item 编写使用说明和硬件接线图
\end{enumerate}

\subsection*{注意点}
\begin{itemize}
    \item ESP32 录音 3s + ASR + LLM 总延迟约 3--5 秒，注意用户体验。
    \item 示波器亮度低时文字可能看不清，适当调节。
    \item 联调时先确保 UART 通信正常，再用示波器查看。
\end{itemize}

% ==========================================
% 六、参考信息
% ==========================================
\chapter{参考信息}

\section{文件清单}

\begin{center}
\renewcommand{\arraystretch}{1.3}
\begin{tabularx}{\textwidth}{|c|X|X|c|}
    \hline
    \heiti 文件 & \heiti 功能 & \heiti 阶段 & \heiti 状态\\
    \hline
    \texttt{AI\_Chat/src/main.cpp} & 主程序（录音、ASR、LLM） & 一 & \taskdone 已完成 \\
    \hline
    \texttt{AI\_Chat/include/config.h} & 配置（WiFi、API、引脚） & 一 & \taskdone 已完成 \\
    \hline
    \texttt{AI\_Chat/lib/microphone/src/microphone.cpp} & INMP441 I2S 驱动（标准 I2S RX） & 一 & \taskdone 已完成 \\
    \hline
    \texttt{AI\_Chat/lib/microphone/src/microphone.h} & INMP441 驱动头文件 & 一 & \taskdone 已完成 \\
    \hline
    \texttt{AI\_Chat/lib/inmp441/src/inmp441.cpp} & INMP441 PDM 驱动（\textbf{已废弃，勿用}） & 一 & \taskfail 废弃 \\
    \hline
    \texttt{AI\_Chat/lib/inmp441/src/inmp441.h} & PDM 驱动头文件（\textbf{已废弃，勿用}） & 一 & \taskfail 废弃 \\
    \hline
    \texttt{AI\_Chat/platformio.ini} & PlatformIO 项目配置 & 一 & \taskdone 已完成 \\
    \hline
    \texttt{AI\_Chat/项目文档/project.tex} & 本开发说明文档 & 全 & \taskdone 已完成 \\
    \hline
    \texttt{EDC\_407/...}（待创建） & STM32 工程（Keil MDK 项目） & 二 & \tasktodo 待创建 \\
    \hline
\end{tabularx}
\end{center}

\section{API 凭据}

\begin{center}
\renewcommand{\arraystretch}{1.3}
\begin{tabularx}{\textwidth}{|c|X|c|}
    \hline
    \heiti 服务 & \heiti 凭据 & \heiti 说明\\
    \hline
    百度 ASR & AppID: 7825090 & 已在 token 中缓存 \\
    \hline
    百度 ASR & API Key: LVZcNhtn584JqOPB9UCsBE4H & — \\
    \hline
    百度 ASR & Secret Key: QpoVPPhETyWK7yLC8L2TLRzrSBALASG3 & — \\
    \hline
    DeepSeek & API Key: sk-717fd68ac7964fcabd7733cf8917f5a8 & 余额请自行检查 \\
    \hline
\end{tabularx}
\end{center}

\section{已知风险}

\begin{enumerate}[leftmargin=2em]
    \item \textbf{PSRAM}：部分 N8R2 板子 PSRAM 可能未焊接。若始终不可用，录音必须 $\le$ 3 秒。
    \item \textbf{API 费用}：百度 ASR 有免费额度，DeepSeek 按 token 计费。
    \item \textbf{模型变更}：DeepSeek 模型名可能更新（\texttt{deepseek-chat} 2026/07/24 废弃），关注官方公告。
    \item \textbf{示波器闪烁}：矢量扫描需调高 XY 模式余晖时间。
    \item \textbf{波特率一致}：ESP32 与 STM32 UART 必须同为 115200。
    \item \textbf{上传端口抖动}：ESP32-S3 上电后 USB-Serial/JTAG 端口可能变化。如上传失败，按 RST 键重试或重新插拔 USB。
    \item \textbf{旧驱动混淆}：\texttt{lib/inmp441/} 目录包含废弃的 PDM 驱动，编辑时注意不要误引。
\end{enumerate}

\end{document}
